%% Section 06 — Appendices
\chapter*{Appendices}
\addcontentsline{toc}{chapter}{Appendices}

\section*{Appendix A: Research Integrity Online}
\addcontentsline{toc}{section}{Appendix A: Research Integrity Online}

[Insert RIO completion certificate. I am on it, next I will add it to this document before submission]

\section*{Appendix B: Coursework}
\addcontentsline{toc}{section}{Appendix B: Coursework}

AIRS / IFN006 is currently underway, this course will end on july, 2026.

\section*{Appendix C: Notation}
\addcontentsline{toc}{section}{Appendix C: Notation}

\paragraph{Indices and sets.}
\begin{tabular}{@{}ll@{}}
$c, i$        & countries (Consultative Parties), $c \in \{1,\dots,C\}$ \\
$k$           & topic, $k \in K$; $|K|$ is the number of topics \\
$y$           & year \\
$d$           & document (Working Paper) \\
$p$           & passage (chunk) of a document \\
$\ell$        & ordinal stance class, $\ell \in \{1,\dots,5\}$ \\
$D(c,y)$      & Working Papers co-sponsored by country $c$ in year $y$ \\
$K(d)$        & topics assigned to document $d$ \\
$P_d$         & passages of document $d$, $P_d=\{p_1,\dots,p_{n_d}\}$ \\
$P_{d,k}$     & passages of $d$ selected as relevant to topic $k$ \\
$A_d$         & co-sponsors of document $d$ \\
\end{tabular}

\paragraph{Embedding space.}
\begin{tabular}{@{}ll@{}}
$m$                  & embedding dimensionality (full Stella space) \\
$\mathbf{e}_p$       & embedding of passage $p$, $\mathbf{e}_p \in \mathbb{R}^{m}$ \\
$\boldsymbol{\mu}_k$ & centroid of topic $k$, $\boldsymbol{\mu}_k \in \mathbb{R}^{m}$ \\
$\rho_{p,k}$         & passage--topic relevance, $\rho_{p,k}=\cos(\mathbf{e}_p,\boldsymbol{\mu}_k)$ \\
$\tau$               & relevance selection threshold \\
\end{tabular}

\paragraph{Stance and hypotheses.}
\begin{tabular}{@{}ll@{}}
$h_k^{+}, h_k^{0}, h_k^{-}$ & for-, neutral-, against-pole hypotheses for topic $k$ \\
$\lambda(d,h)$  & null-premise-debiased entailment log-ratio (DCPMI) \\
$\tilde\pi_{d,k}$ & three-pole distribution (stance classifier) \\
$\pi_{d,k}$     & five-class ordinal distribution, $\pi_{d,k}\in\Delta_5$ \\
$\Delta_5$      & probability simplex, $\{\pi\in\mathbb{R}^5 : \pi_\ell\ge 0,\ \sum_\ell \pi_\ell = 1\}$ \\
$v$             & class values, $v=(-1,-\tfrac12,0,\tfrac12,1)$ \\
$s_{d,k}$       & signed stance of $d$ on topic $k$, $s_{d,k}\in[-1,1]$ \\
\end{tabular}

\paragraph{Aggregation and tensors.}
\begin{tabular}{@{}ll@{}}
$w_{d,c}$           & attribution weight: $1$ (full credit) or $1/|A_d|$ (fractional) \\
$S_{c,k,y}$         & support of country $c$ on topic $k$ in year $y$, Eq.~(\ref{eq:aggregate}) \\
$S_{\cdot,\cdot,y}$ & year-$y$ stance slice, $S_{\cdot,\cdot,y}\in\mathbb{R}^{C\times|K|}$ \\
$O_{c,k,y}$         & opposition of country $c$ on topic $k$ in year $y$ \\
\end{tabular}

For these tensors, high sparsity is expected because absence of evidence is recorded as missing, not as zero.

\paragraph{Stance classifier (Chapter~1).}
For topic $k$, let $h_k^{+}, h_k^{0}, h_k^{-}$ be the for-, neutral-, and
against-pole hypotheses. For passage set $d$ and hypothesis $h$, the debiased evidence is the null-premise-corrected log-ratio,
\[
  \lambda(d,h) = \operatorname{clip}\!\Big(
     \log P(\text{entail}\mid d,h) - \log P(\text{entail}\mid \varnothing,h),
     \ \pm 5 \Big),
\]
this removes each hypothesis's marginal acceptability. The three-pole distribution
combines entailment of each pole with contradiction of its opposite,
\[
  \tilde\pi_{d,k} = \operatorname{softmax}\!\begin{pmatrix}
     \lambda_{\mathrm{ent}}(h_k^{+}) + \lambda_{\mathrm{con}}(h_k^{-}) \\[2pt]
     \lambda_{\mathrm{ent}}(h_k^{0}) \\[2pt]
     \lambda_{\mathrm{ent}}(h_k^{-}) + \lambda_{\mathrm{con}}(h_k^{+})
  \end{pmatrix},
\]
these values are expanded to the five ordinal classes $\pi_{d,k}$ and reduced to the
scalar stance by the expected value of Equation~\ref{eq:stance}.

\section*{Appendix D: Opposition-layer workflow}
\addcontentsline{toc}{section}{Appendix D: Opposition-layer workflow}

\begin{figure}[htbp]
  \centering
  \resizebox{\textwidth}{!}{% Chapter 2 — Opposition Layer (idea sketch, TikZ).
% Requires in preamble: \usepackage{tikz}, amsmath, cleveref, xcolor[dvipsnames],
% \usetikzlibrary{shapes.geometric, arrows.meta, positioning, calc, fit, backgrounds}
\begin{tikzpicture}[
    font=\small,
    node distance = 7mm and 12mm,
    >={Stealth[length=2.4mm]},
    every node/.style = {align=center},
    proc/.style    = {rectangle, rounded corners=1pt, draw=black!60,
                       fill=NavyBlue!8, text width=54mm, inner sep=4pt},
    term/.style    = {rectangle, rounded corners=8pt, draw=black!60,
                      fill=black!8, text width=54mm, inner sep=4pt},
    note/.style    = {rectangle, rounded corners=1pt, draw=black!30, dashed,
                      fill=Goldenrod!12, text width=46mm, align=left,
                      font=\footnotesize, inner sep=4pt},
    arr/.style     = {->, draw=black!70},
    merge/.style   = {-, draw=black!70},
    link/.style    = {-, draw=black!35, dashed}
  ]

  % ── main spine ───────────────────────────────────────────────
  \node[term]                       (start)   {ATCM Final Reports (text input)};
  \node[proc, below=of start]       (extract) {Extract veto statements\\ $\to$ (country, target, polarity) triples\\ + rationale \& source quote};
  \node[proc, below=of extract]     (map)     {Map targets to Chapter~1 topics $K$};

  % ── parallel split (primary / secondary) ─────────────────────
  \coordinate (fork) at ($(map.south)+(0,-7mm)$);
  \node[proc] (primary)   at ($(map.south)+(-32mm,-20mm)$) {\textbf{Primary:} opposition tensor $O$\\ (country $\times$ topic $\times$ year)};
  \node[proc] (secondary) at ($(map.south)+( 32mm,-20mm)$) {\textbf{Secondary (exploratory):} country propensity to reject proposals per ATCM};

  \node[proc, below=42mm of map]    (net)     {Signed preference network\\ = tensor pair $(S,O)$};
  \coordinate (join) at ($(net.north)+(0,8mm)$);

  \node[proc, below=14mm of net]    (val)     {Test $(S,O)$ for predictive content against observed ATCM outcomes + \textcite{Gardiner2025}};
  \node[proc, below=14mm of val]    (toch3)   {$(S,O)\ \to\ $ input to Chapter~3};
  \node[term, below=of toch3]       (stop)    {End};

  % ── side notes ───────────────────────────────────────────────
  \node[note, right=of extract] (n1) {Attribution harder than the support layer (sparse text).};
  \node[note, left=of primary]  (n2) {Veto signal \emph{against} a proposal --- not the inverse of support $S$.};
  \node[note, right=of net]     (n3) {$S$ from Chapter~1.};
  \node[note, right=of val]     (n4) {Independent research question; risk: exogenous geopolitics (e.g.\ 2022 Ukraine breakdown) not captured in institutional text.};

  % ── forward edges ────────────────────────────────────────────
  \draw[arr] (start)   -- (extract);
  \draw[arr] (extract) -- (map);
  \draw[merge] (map)   -- (fork);
  \draw[arr] (fork) -| (primary.north);
  \draw[arr] (fork) -| (secondary.north);
  \draw[merge] (primary.south)   |- (join);
  \draw[merge] (secondary.south) |- (join);
  \draw[arr] (join)  -- (net);
  \draw[arr] (net)   -- (val);
  \draw[arr] (val)   -- (toch3);
  \draw[arr] (toch3) -- (stop);

  % ── note links ───────────────────────────────────────────────
  \draw[link] (extract.east) -- (n1.west);
  \draw[link] (primary.west) -- (n2.east);
  \draw[link] (net.east)     -- (n3.west);
  \draw[link] (val.east)     -- (n4.west);

  % ── partition background ─────────────────────────────────────
  \begin{scope}[on background layer]
    \node[draw=Maroon!40, fill=Maroon!4, rounded corners,
          fit=(val), inner sep=6pt,
          label={[Maroon!70!black, font=\footnotesize\itshape]above:External validation}] {};
  \end{scope}

\end{tikzpicture}
}
  \caption{Chapter~2 opposition layer sketch: from text data to opposition tensor~$O$, forming the signed preference network~$(S,O)$.}
  \label{fig:ch2-workflow}
\end{figure}

\section*{Appendix E: Extended literature notes}
\addcontentsline{toc}{section}{Appendix E: Extended literature notes}

\paragraph{Wordfish \parencite{Slapin2008}.}
Its aim was to estimate party positions over time without hand-coded reference texts. The model
treats each word count as a Poisson draw whose expected frequency depends on the party's position
on a single latent dimension, and it recovered credible left--right trajectories for German
parties from their election manifestos (1990--2005).

\paragraph{Wordshoal \parencite{Lauderdale2016}.}
Their model scales each debate separately with Wordfish and then combines the debate-level
positions into a general dimension through a Bayesian factor analysis. Applied to speeches in the
Irish D\'ail and the US Senate, it recovered within-party disagreement that voting records hide.
Each debate scale is whatever axis dominates that debate's word frequencies, which need not be
the substantive for-or-against axis of interest. Lauderdale and Herzog get away with this because
legislative debates are adversarial, so the dominant axis is usually government versus
opposition; ATCM Working Papers have no such structure.

\paragraph{Hedging in diplomatic text.}
Hedging means cautious, non-committal phrasing: a party rarely writes ``we oppose this measure'',
it writes ``further consideration may be warranted''. Positions on opposite sides also share the
same vocabulary of safeguards: a phrase such as ``subject to appropriate environmental
safeguards'' appears in permissive and restrictive papers alike. Example pole hypotheses for a
topic: $h_k^{+}$ ``tourism should be further regulated'', $h_k^{-}$ ``tourism regulation should
be relaxed''.

\paragraph{The UN General Debate corpus \parencite{Baturo2017}.}
Their corpus collects over 7300 annual General Debate statements (1970--2014), one per country
per year. They derive preferences in two ways, both from word frequencies: Wordscores, a
supervised scaling in which the analyst defines the axis in advance through reference texts (US
and Russian statements anchor a US--Russia dimension), and correspondence analysis, an
unsupervised reduction whose dimensions are interpreted afterwards.
